Redefining the International System of Units (SI) second requires
independent frequency-ratio measurements between multiple optical
clocks in different laboratories, with an overall measurement
uncertainty below \refThreshold{} and consistent results. However,
although several optical clocks have reported self-evaluated
uncertainties below \selfEvalBound, validation of these self-assessed
values through frequency-ratio comparisons has never been achieved. In
this study, we compared a ytterbium optical lattice clock in Wuhan with
a strontium optical lattice clock in Shanghai, approximately
$\siteKm$\,km away, using a phase-stabilised fibre link spanning
$\fibreKm$\,km. Based on data collected over \nDays~consecutive days
($\sim$\,\nHours~hours), we measured the Yb/Sr frequency ratio as
Yb/Sr~=~\Rheadline(\RheadlineUnc), with a Bayesian statistical
uncertainty of \uStatShort$\times10^{-18}$ and a total uncertainty of
\uTotalShort$\times10^{-18}$. The value differs from that of
NIST--JILA~\refNistV(\refNistU) by \dVsNistE{}, within the threshold
range required for redefining the SI second. To our knowledge, this is
the first inter-city, cross-species frequency-ratio measurement to meet
the uncertainty requirements outlined in the roadmap for redefining the
SI second. Furthermore, together with the NIST--JILA frequency-ratio
result, this achievement satisfies the mandatory criteria for changing
the definition of the second, confirming the validity of uncertainty
assessments for both ytterbium and strontium atomic clocks.
